\documentclass[12pt]{beamer}
\usetheme{Berkeley}
\usepackage[utf8]{inputenc}
\usepackage[german]{babel}
\usepackage[T1]{fontenc}
\usepackage{amsmath}
\usepackage{amsfonts}
\usepackage{amssymb}
\usepackage{graphicx}
\author{Nicas Heydorn \and Kevin Kahrau \and Christofer Ostwald}
\title{Deepnet \& Darknet}
%\setbeamercovered{transparent} 
%\setbeamertemplate{navigation symbols}{}
\logo{\includegraphics[height=1.6cm]{assets/0pVtHi3d.jpg}}
\institute{Nordakademie - Hochschule der Wirtschaft} 
%\date{} 
\subject{Technische Grundlagen 2} 
\begin{document}

\begin{frame}
\titlepage
\end{frame}

\logo{\includegraphics[width=1.8cm]{assets/NAK-Productions-Logo.png}}
\begin{frame}
\tableofcontents
\end{frame}


\section{Definition}

\logo{\includegraphics[height=1.6cm]{assets/nak_letters_red_0713.png}}
\begin{frame}{Definition}
	Deep Web - \emph{Verstecktes Web}
	\begin{itemize}
		\item Teil des WWW welches nicht über Suchmaschiene auffindbar ist
		\item Gegensatz Visible Web oder Surface Web
\item Inhalte sind nicht frei zugänglich oder nicht durch Suchmaschienen indezierbar
	\end{itemize}
\end{frame}



\section{Arten}

\logo{\includegraphics[height=1.6cm]{assets/1404213_1_high.png}}
\begin{frame}{Opaque Web}
	\begin{itemize}
		\item könnten technisch indexiert werden ist jedoch zu aufwendig für den erwarteten Nutzen
		\item Webcrawler berücksichten oft nur 5-6 Verzeichnisebene und finden somit keine tief verschachtelten Dokumente
		\item Nur teilweise indezierbare Dateiformate (z.B PDF)
	\end{itemize}
\end{frame}

\logo{\includegraphics[height=1.6cm]{assets/nak.png}}
\begin{frame}{Private Web}
	\begin{itemize}
		\item keine Indexierung durch Zugriffsbeschränkung des Webmasters 
		\begin{itemize}
			\item Intranet
			\item Authentication
			\item IP-Whitelist
			\item Meta-Tag \textit{noindex/nofollow/noimageindex}
			\item \glqq Robots Exclusion Standard\grqq
		\end{itemize}	
	\end{itemize}
\end{frame}

\logo{\includegraphics[height=1.6cm]{assets/flytjenesten_logo_1.png}}
\begin{frame}{Proprietary Web}
	\begin{itemize}
		\item Zustimmung von Nutzungsbedingen erforderlich
		\item Passwort erforderlich
		\item z.B. webbasierte Fachdatenbanken
	\end{itemize}
\end{frame}

\logo{\includegraphics[width=1.6cm]{assets/my-hair-nak-logo.jpg}}
\begin{frame}{Invisible Web}
	\begin{itemize}
		\item Webseiten welche technisch indexiert werden könnten, jedoch aus wirtschaftlichen oder strategischen Gründen nicht indexiert werden
		\item Beispiel: Datenbanken mit Webformular
	\end{itemize}
\end{frame}

\logo{\includegraphics[height=1.6cm]{assets/Original_Nak_Logo_full.png}}
\begin{frame}{Truly Invisible Web}
	\begin{itemize}
		\item Webseiten welche technisch (noch) nicht indexiert werden können
		\item Beispiele:
		\begin{itemize}
			\item alte Datenbankformate
			\item Dokumente welche nicht im Browser angezeigt werden können
			\item Nicht-Standardformate
			\item komplexe Dateiformate (Grafikformate)
			\item Komprimierte Daten
			\item Webseiten die nur über Benutzernavigation zu bedienen sind
		\end{itemize}
	\end{itemize} 
\end{frame}

\end{document}